\documentclass{beamer}

% --- Use the following to create a printout for the presentation ---
% \documentclass[handout]{beamer}

% \usepackage{pgfpages}
% \pgfpagesuselayout{2 on 1}[a4paper,border shrink=5mm]

\setbeamertemplate{theorems}[numbered]
% \beamerdefaultoverlayspecification{<+->}

\usetheme{gotham}

	\gothamset{
		% background=dark,
		numbering= framenumber,
		% tocframe template= gotham simple,
		parttocframe default=off,
		sectiontocframe default=off,
		subsectiontocframe default=off,
    progressbar style=rounded box
	}

\usepackage{standalone}
\usepackage{changepage}
\usepackage{algorithm}
\usepackage{algpseudocode}


\usepackage[backend=biber, style=numeric, sorting=none]{biblatex}
\addbibresource{cs598.bib}

\usepackage{appendixnumberbeamer}
\gdef\appendixtotalframenumber{1}

\usepackage{tabularx}
\usepackage{booktabs}
\usepackage{physics2}
\usepackage{derivative}
\usepackage[noabbrev,capitalize]{cleveref}
\usepackage{nicefrac}
\usepackage{todonotes}

\usepackage[acronym,toc]{glossaries}
\include{src/acronymns}
\makeglossaries

\title[]{Parallel Matrix-Free Spectral Elements for Solving 3D Poisson Equation}
\subtitle[]{CS 598 Midterm}
\date[]{November 2026}
\author[]{\inst{1}Liam Pohlmann, \inst{2}Alex Shaffer}
\institute{\inst{1} Department of Nuclear, Plasma, and Radiological Engineering \and \inst{2} Department of Computer Science }

\begin{document}

\maketitle

\begin{frame}[toc]{Table of Contents}%
	\small
	\tableofcontents%[hideallsubsections]
\end{frame}

\section{Project Description}
\begin{frame}{Toy Problem}
  For this implementation, we look to solve the following toy problem:

  Find $\tilde{u}\in \mathcal{R}^3$ s.t.
  \begin{equation}
    \label{eq:poisson}
    -\nabla^2 \tilde{u}=f(x,y,z),\quad \tilde{u}\in \Omega,\quad u\in\partial \Omega=0.
  \end{equation}
  for $x,y,z\in [0,1]^3$.
\end{frame}

\section{Numerical Model}
\subsection{Weak Form}
\begin{frame}{Weak Form of PDE}
  To develop the weak form, we first approximate the exact solution with a numerical solution, $\tilde{u}\approx u$, then multiply both sides of \cref{eq:poisson} by a test function, $v$, and integrate over the domain:
  \begin{equation}
    -\int v\nabla^2 u\odif{V} = \int vf\odif{V}
  \end{equation}
  Using vector identities and integration by parts leaves:
  \begin{equation}
    \label{eq:weak-form}
    -\int \nabla u\cdot\nabla v\odif{V}+\oint v\nabla u\cdot\hat{n}\odif{S}=\int fv\odif{V}.
  \end{equation}
\end{frame}

\begin{frame}{Galerkin Approximation}
  We first choose the approximation $u$ to be a sum of Lagrange basis functions:
  \begin{equation}
    \label{eq:lagrange-expansion}
    u(x,y,z)=\sum_{k=0}^{N_z}\sum_{j=0}^{N_y}\sum_{i=0}^{N_x}u_{ijk}\ell_i(x)\ell_j(y)\ell_k(z)
  \end{equation}
  for $N_x$, $N_y$, and $N_z$ elements in the $x$, $y$, $z$ directions, respectively.
  For a Galerkin projection, we choose:
  \begin{equation}
    \label{eq:test-galerkin}
    v(x,y,z)=\sum_{n=0}^{N_z}\sum_{m=0}^{N_y}\sum_{l=0}^{N_x}v_{lmn}\ell_l(x)\ell_m(y)\ell_n(z).
  \end{equation}
\end{frame}

\begin{frame}{Galerkin Projection}
  Using \cref{eq:lagrange-expansion} and \cref{eq:test-galerkin} in \cref{eq:weak-form} gives a set of $(N_x+1)(N_y+1)(N_z+1)$ equations.
    [UNFINISHED]
\end{frame}

\subsection{Spectral Element Method}
\begin{frame}{Spectral Element Method}
  The defining feature of the spectral element method (SEM) is the selection of internal element nodes corresponding to the Gauss-Lobatto-Legendre (GLL) rule.
  From \cite{devilleHighorderMethodsIncompressible2002}:
  \begin{equation}
    \label{eq:gll}
    \begin{aligned}
       & \int_{-1}^1 u(\xi)\odif{\xi} & \approx & \sum_{k=0}^{N}\rho_k u\left(\xi_k\right)                                  \\
       & \xi_k                        & =       & \begin{cases}
                                                    -1,                        & k=0              \\
                                                    \text{zeros of }L'_N(\xi), & 1\leq k \leq N-1 \\
                                                    1,                         & k=N              \\
                                                  \end{cases}                           \\
       & \rho_k                       & =       & \frac{2}{N\left(N+1\right)}\frac{1}{\left[L_N\left(\xi_k\right)\right]^2}
    \end{aligned}
  \end{equation}
  The above rule is exact for polynomials degree up to $2N-1$.
\end{frame}

\begin{frame}{Spectral Element Method}
  Using the points in \cref{eq:gll}, which we approximate the integrals for the stiffness and mass matrices.
  Using \textit{generic} indices $i,j$, we have:
  \begin{subequations}
    \begin{equation}
      \left(\hat{M}_n\right)_{ij}  \approx \frac{L_n}{2}\sum_{k=0}^{N}\rho_k\ell_i (\mu_k)\ell_j (\mu_k)=\begin{cases}
        \rho_i, & i=j     \\
        0,      & i\neq j
      \end{cases}
    \end{equation}
    \begin{equation}
      \left(\hat{K}_n\right)_{ij}  \approx \frac{2}{L_n}\sum_{k=0}^{N}\rho_k \odv{\ell_i}{\mu}(\mu_k)\odv{\ell_j}{\mu}(\mu_k)
    \end{equation}
  \end{subequations}
  Notice that because the GLL quadrature points are also the nodes for the Lagrange interpolants, the mass matrix is \emph{diagonal}.
  The stiffness matrix, however, is full.
\end{frame}

\section{Implementation}

\subsection{Kronecker Product Evaluation}
\begin{frame}{Kronecker Product}
  For three square matrices of the same size, $A, B, C\in \mathbb{R}^{(N+1) \times (N+1)}$, the Kronecker product can be expanded using $(N+1)\times (N+1)$ identity matrices:
  \begin{equation}
    A\otimes B\otimes C=\left(A\otimes I\otimes I\right)\left(I\otimes B\otimes I\right)\left(I\otimes I\otimes C\right).
  \end{equation}
  We will be interested in finding the product, $\left(A\otimes I\otimes I\right)\left(I\otimes B\otimes I\right)\left(I\otimes I\otimes C\right)\mathbf{u}$, where $\mathbf{u}$ is a length $(N+1)^3$ vector.
  Let:
  \begin{subequations}
    \begin{equation}
      \mathbf{w}  =\left(I\otimes I\otimes C\right)\mathbf{u}
    \end{equation}
    \begin{equation}
      \mathbf{v}  =\left(I\otimes B\otimes I\right)\mathbf{w}
    \end{equation}
    \begin{equation}
      \mathbf{y}  =\left(A\otimes I\otimes I\right)\mathbf{v}
    \end{equation}
  \end{subequations}
\end{frame}

\begin{frame}{Kronecker Product}
  Define the flat indexing:
  \begin{equation}
    \hat{\imath} (i,j,k)=i+j(N+1)+k(N+1)^2.
  \end{equation}
  Then we have the following definitions:
  \begin{subequations}
    \begin{equation}
      w_{\hat{\imath}}= w_{ijk}=\sum_{p=0}^{N}C_{ip}u_{pjk}
    \end{equation}
    \begin{equation}
      v_{\hat{\imath}}=v_{ijk}=\sum_{p=0}^{N}B_{jp}w_{ipk}
    \end{equation}
    \begin{equation}
      y_{\hat{\imath}}=y_{ijk}=\sum_{p=0}^{N}A_{kp}v_{ijp}.
    \end{equation}
  \end{subequations}
  Note that in doing this, we reduce the matrix-vector product from $\mathcal{O}\left((N+1)^6\right)$ to $\mathcal{O}\left((N+1)^4\right)$, recovering two orders of magnitude.
\end{frame}

\subsection{Gather/Scatter Operator}
\begin{frame}{Gather-Scatter Operator}

\end{frame}

\subsection{Iterative Solver}
\begin{frame}{Conjugate Gradient}
    Conjugate Gradient (CG) solves the linear system $A\mathbf{x}=\mathbf{b}$ for a symmetric positive definite matrix $A$ iteratively.

    It is very well suited for SEM because it:
    \begin{itemize}
        \item Only requires the action of the matrix $A$ on a vector.
        \item Only requires a couple global reductions per iteration.
    \end{itemize}

\end{frame}

\begin{frame}{Conjugate Gradient}
    \begin{algorithm}[H]
      \caption{Unpreconditioned CG for $A u = b$}
      \begin{algorithmic}[1]
        \State $u \gets 0$, \quad $r \gets b$, \quad $p \gets r$
        \State $\rho \gets r^{T} r$
        \While{$\sqrt{\rho} \ge \varepsilon \, \lVert b \rVert$}
          \State $w \gets A p$ \Comment{Kronecker products + gather-scatter}
          \State $\alpha \gets \rho / (p^{T} w)$ \Comment{\texttt{MPI\_Allreduce}}
          \State $u \gets u + \alpha p$
          \State $r \gets r - \alpha w$
          \State $\rho_{\text{new}} \gets r^{T} r$ \Comment{\texttt{MPI\_Allreduce}}
          \State $p \gets r + (\rho_{\text{new}} / \rho) \, p$
          \State $\rho \gets \rho_{\text{new}}$
        \EndWhile
      \end{algorithmic}
    \end{algorithm}
\end{frame}


\section{Results}
\begin{frame}{$p$-type Convergence}

\end{frame}

\begin{frame}{$h$-type Convergence}

\end{frame}

\begin{frame}{Strong Scaling}

\end{frame}

\begin{frame}{Weak Scaling}

\end{frame}

\begin{frame}{Additional Notes}
  \begin{itemize}
    \item Most of the compute time is eaten up by Kronecker product evaluations.
    \item The biggest speedup was by telling our Kronecker product routine ahead of time whether or not a matrix was diagonal.
  \end{itemize}
\end{frame}

\section{References}
\begin{frame}[noframenumbering,plain,allowframebreaks]
	\frametitle{References}
	\printbibliography[heading=none]
\end{frame}


\end{document}
